%% $Id: preamble.tex 20694 2010-09-21 18:47:18Z menski $

%% $Id: preamble-commands.tex 20694 2010-09-21 18:47:18Z menski $

\usepackage{xspace}
\usepackage{ifthen}
\usepackage{ifpdf}

\makeatletter

\providecommand{\IfPackageLoaded}[2]{\@ifpackageloaded{#1}{#2}{}}
\providecommand{\IfPackageNotLoaded}[2]{\@ifpackageloaded{#1}{}{#2}}
\providecommand{\IfElsePackageLoaded}[3]{\@ifpackageloaded{#1}{#2}{#3}}

\def \mail {}
\newcommand{\email}[1]{\def \mail {#1}}

\@ifundefined{frontmatter}{%
   \newcommand{\frontmatter}{%
      %In Roemischen Buchstaben nummerieren (i, ii, iii)
      \pagenumbering{roman}
   }
}{}
\@ifundefined{mainmatter}{%
   % scrpage2 benoetigt den folgenden switch
   % wenn \mainmatter definiert ist.
   \newif\if@mainmatter\@mainmattertrue
   \newcommand{\mainmatter}{%
      % -- Seitennummerierung auf Arabische Zahlen zuruecksetzen (1,2,3)
      \pagestyle{fancy}%%
      \pagenumbering{arabic}%%
   }
}{}
\@ifundefined{backmatter}{%
   \newcommand{\backmatter}{
      %In Roemischen Buchstaben nummerieren (i, ii, iii)
      \pagenumbering{Roman}
   }
}{}

% \def\thickhrulefill{\leavevmode \leaders \hrule height 1pt\hfill \kern \z@}
% \renewcommand{\maketitle}{\begin{titlepage}%
%     \let\footnotesize\small
%     \let\footnoterule\relax
    % \parindent \z@
    % \reset@font
    % \null
    % \vskip 10\p@
    % \hbox{\mbox{%
    %     \hspace{4pt}%
      %   		\includegraphics[width=7em]{images/uni-logo}%
      %   \hspace{4pt}
      %   }%
      % \vrule depth 0.80\textheight%
      % \mbox{\hspace{2em}}
      % \vtop{% %%%%%%%%%%%%%%%%%%
      %   \vskip 20\p@
        % \begin{flushleft}
        %   \Huge \bfseries \@title \par
        % \end{flushleft}
        % \vskip 60\p@
        % \begin{flushleft}
        %   \LARGE BACHELORARBEIT \par 
        %   \vskip 5\p@
        %   \large zur Erlangung des akademischen Grades\\\Large\textit{Bachelor of Science}\\\large im Studiengang Informatik \par
        %   \vskip 5\p@
        %   \LARGE \@author \par
        %   \vskip 1\p@
        %   \large \mail \par
        % \end{flushleft}
	% 	\vskip 90\p@
        % \begin{flushleft}
        %   \large \@publishers \par
        % \end{flushleft}
        % \vskip 30\p@
        % \begin{flushleft}
        %   \large \@date \par
    %     \end{flushleft}
    %     \vfil
    %     }}        
    % \null
%   \end{titlepage}%
%   \setcounter{footnote}{0}%
% }

%\def\thickhrulefill{\leavevmode \leaders \hrule height 1pt\hfill \kern \z@}
%\renewcommand{\maketitle}{\begin{titlepage}%
%    \let\footnotesize\small
%    \let\footnoterule\relax
%    \parindent \z@
%    \reset@font
%    \null
%    \vskip 10\p@
%    \hbox{\mbox{%
%        \hspace{4pt}%
%			\includegraphics[width=7em]{images/uni-logo}%
%        \hspace{4pt}
%        }%
%      \vrule depth 0.80\textheight%
%      \mbox{\hspace{2em}}
%      \vtop{% %%%%%%%%%%%%%%%%%%
%        \vskip 20\p@
%        \begin{flushleft}
%          \Huge \bfseries \@title \par
%          \LARGE \bfseries \@subtitle \par
%        \end{flushleft}
%        \vskip 60\p@
%        \begin{flushleft}
%          \LARGE Bachelor Thesis \par 
%          \vskip 5\p@
%          \large by \par
%          \vskip 5\p@
%          \LARGE \@author \par
%          \vskip 1\p@
%          \large \mail \par
%        \end{flushleft}
%		\vskip 120\p@
%        \begin{flushleft}
%          \large \@publishers \par
%        \end{flushleft}
%        \vskip 30\p@
%        \begin{flushleft}
%          \large \@date \par
%        \end{flushleft}
%        \vfil
%        }}        
%    \null
%  \end{titlepage}%
%  \setcounter{footnote}{0}%
%}

\makeatother

% Encoding
\usepackage[T1]{fontenc}
\usepackage[%
%	latin1,
	utf8
]{inputenc}

% Layout
\usepackage{fancyhdr}

% Refernces
\usepackage[]{prettyref}
%\newrefformat{cha}{Kapitel \ref{#1}}
%\newrefformat{sec}{Abschnitt \ref{#1}}
%\newrefformat{tab}{Tabelle \ref{#1} auf Seite \pageref{#1}}
%\newrefformat{fig}{Abbildung \ref{#1} auf Seite \pageref{#1}}
%\newrefformat{fig}{Abbildung \ref{#1}}
\newrefformat{fig}{Figure \ref{#1}}
%\newrefformat{lst}{Listing \ref{#1} auf Seite \pageref{#1}}

\usepackage{todonotes}

% Grafik
% \usepackage[%
% 	table
% ]{xcolor}
\usepackage{graphicx}
\usepackage{epstopdf}

% Mathe
\usepackage[%
	centertags,
	sumlimits,
	intlimits,
	namelimits,
	fleqn
]{amsmath}
\usepackage[
	all,
	warning
]{onlyamsmath}
\usepackage{fixmath}

\usepackage{tikz}
\usetikzlibrary{arrows}

% Berechnungen
\usepackage{calc}

% Textauszeichnung
\usepackage{soul}
\usepackage[normalem]{ulem}
\usepackage{url}
\usepackage[%
	babel,
	german=quotes,
	english=british
]{csquotes}
\usepackage{ragged2e}
\usepackage{multicol}
\usepackage[
	colorlinks=false,
	pdfborder={0 0 0},
]{hyperref}

% Floats
\usepackage{float}
\usepackage{flafter}
\usepackage[%
	section,
	above
]{placeins}
\usepackage[lofdepth=2]{subfig}
\usepackage{wrapfig}
\setlength{\intextsep}{0.5\baselineskip}

% Tabellen
\usepackage{booktabs}
\usepackage{multirow}
\usepackage{dcolumn}
\usepackage{ltxtable}

% Wissenschaft
\usepackage{units}

% Listings
%\usepackage{algorithmic}
%\usepackage{algorithm}
\usepackage[ruled,vlined,linesnumbered]{algorithm2e}
\usepackage{listings}
\lstset{ %
language=Java,                % choose the language of the code
basicstyle=\footnotesize,       % the size of the fonts that are used for the code
numbers=left,                   % where to put the line-numbers
numberstyle=\footnotesize,      % the size of the fonts that are used for the line-numbers
stepnumber=1,   	            % the step between two line-numbers. If it's 1 each line 
                                % will be numbered
numbersep=5pt,                  % how far the line-numbers are from the code
backgroundcolor=\color{white},  % choose the background color. You must add \usepackage{color}
showspaces=false,               % show spaces adding particular underscores
showstringspaces=false,         % underline spaces within strings
showtabs=false,                 % show tabs within strings adding particular underscores
%frame=single,	                % adds a frame around the code
tabsize=2,	                % sets default tabsize to 2 spaces
captionpos=b,                   % sets the caption-position to bottom
breaklines=true,                % sets automatic line breaking
breakatwhitespace=false,        % sets if automatic breaks should only happen at whitespace
%title=\lstname,                 % show the filename of files included with \lstinputlisting;
                                % also try caption instead of title
escapeinside={\%*}{*)},         % if you want to add a comment within your code
morekeywords={*,...}            % if you want to add more keywords to the set
}


%%% Local Variables: 
%%% mode: latex
%%% TeX-master: "../master"
%%% End: 
